\documentclass[../distribution_theory_notes.tex]{subfiles}
\begin{document}
\section{Aula 25 - 22 de Novembro, 2024}
\subsection{Motivações}
\begin{itemize}
	\item Transformada de Fourier de funções em $L^{1}(\mathbb{R}^{n})$;
	\item Propriedades operacionais da Transformada de Fourier;
	\item Lema de Riemann--Lebesgue.
\end{itemize}

\subsection{Transformada de Fourier de Funções $L^{1}(\mathbb{R}^{n})$}
\begin{def*}
	Seja $f\in L^{1}(\mathbb{R}^{n})$. Usaremos a notação
	\[
		x\cdot y=x_{1}y_{1}+\cdots+x_{n}y_{n},
	\]
	para o produto escalar euclidiano, e definimos:
	\begin{itemize}
		\item[(i)] A \textbf{transformada de Fourier de $f$} em $\xi\in\mathbb{R}^{n}$ é o número
		      \[
			      \widehat f(\xi)
			      \coloneqq
			      \int_{\mathbb{R}^{n}}f(x)e^{-2\pi i x\cdot\xi}\,dx.
		      \]

		\item[(ii)] A \textbf{transformada de Fourier inversa de $f$} em $x\in\mathbb{R}^{n}$ é
		      \[
			      \check f(x)
			      \coloneqq
			      \int_{\mathbb{R}^{n}}f(\xi)e^{2\pi i x\cdot\xi}\,d\xi.
		      \]
		      Note que
		      \[
			      \check f(x)=\widehat f(-x).
		      \]

		\item[(iii)] O operador \textbf{transformada de Fourier}, ou apenas \textbf{Transformada de Fourier}, é, por definição, o operador
		      \[
			      \mathcal{F}:L^{1}(\mathbb{R}^{n})\longrightarrow X,
		      \]
		      onde
		      \[
			      X\coloneqq
			      \left\{\widehat f:\mathbb{R}^{n}\longrightarrow\mathbb{C};\
			      \widehat f(\xi),\ \xi\in\mathbb{R}^{n},\
			      \text{para algum }f\in L^{1}\right\},
		      \]
		      dado por
		      \[
			      f\longmapsto \mathcal{F}(f)=\widehat f,
		      \]
		      onde $\widehat f$ é dada por (i).
	\end{itemize}
	A linearidade de $\mathcal{F}$ é imediata a partir da linearidade da integral. \(\square\)
\end{def*}

Comecemos destacando que a Transformada de Fourier é um processo que regulariza as funções.

\begin{prop*}
	Se $f\in L^{1}(\mathbb{R}^{n})$, então
	\[
		\widehat f\in \mathcal{C}_{B}(\mathbb{R}^{n})
		\coloneqq
		\{\text{funções contínuas limitadas}\},
	\]
	com
	\[
		\|\mathcal{F}f\|_{\infty}\leq \|f\|_{L^{1}}.
	\]
	Em particular, $\mathcal{F}$ é contínua.
\end{prop*}

\begin{proof*}
	Com efeito, $\widehat f\in L^{\infty}$, porque
	\[
		|\widehat f(\xi)|
		\leq
		\int_{\mathbb{R}^{n}}|f(x)|\,dx
		=\|f\|_{L^{1}},
	\]
	já que $e^{-2\pi i x\cdot\xi}$ tem valor absoluto $1$. Portanto,
	\[
		\|\widehat f\|_{\infty}\leq\|f\|_{L^{1}},
		\qquad \forall f\in L^{1}.
	\]

	Agora, se $\xi_{j}\to\xi_{0}$ em $\mathbb{R}^{n}$, então
	\[
		f(x)e^{-2\pi i x\cdot\xi_{j}}
		\longrightarrow
		f(x)e^{-2\pi i x\cdot\xi_{0}},
		\qquad \forall x\in\mathbb{R}^{n}.
	\]
	De fato, fixando um representante para $[f]$, isso decorre da continuidade da exponencial complexa. Além disso,
	\[
		\left|f(x)e^{-2\pi i x\cdot\xi_{j}}\right|
		\leq |f(x)|,
		\qquad \forall x\text{ e }j,
	\]
	com $f\in L^{1}$. Pelo Teorema da Convergência Dominada,
	\[
		\widehat f(\xi_{j})
		=
		\int_{\mathbb{R}^{n}}f(x)e^{-2\pi i x\cdot\xi_{j}}\,dx
		\longrightarrow
		\int_{\mathbb{R}^{n}}f(x)e^{-2\pi i x\cdot\xi_{0}}\,dx
		=\widehat f(\xi_{0}),
	\]
	e, portanto, $\widehat f$ é contínua em $\xi_{0}$. \qedsymbol
\end{proof*}

A seguir, são listadas as primeiras propriedades de $\mathcal{F}$ relacionadas a algumas operações que podemos realizar com as funções.

\subsection{Propriedades da Transformada de Fourier}
Sejam $f,g\in L^{1}(\mathbb{R}^{n})$. Valem as propriedades a seguir.

\medskip
\noindent\textbf{(I)} Para todo $h,\xi\in\mathbb{R}^{n}$,
\[
	\widehat{\tau_{h}f}(\xi)
	=e^{-2\pi i h\cdot\xi}\widehat f(\xi).
\]

\begin{proof*}
	Pelo Teorema da Mudança de Variáveis, fazendo $y=x-h$, temos
	\begin{align*}
		\widehat{\tau_{h}f}(\xi)
		 & =\int_{\mathbb{R}^{n}}f(x-h)e^{-2\pi i x\cdot\xi}\,dx   \\
		 & =\int_{\mathbb{R}^{n}}f(y)e^{-2\pi i (y+h)\cdot\xi}\,dy \\
		 & =e^{-2\pi i h\cdot\xi}\widehat f(\xi).
	\end{align*}
	\qedsymbol
\end{proof*}

\medskip
\noindent\textbf{(II)} Para todo $h,\xi\in\mathbb{R}^{n}$,
\[
	\widehat{\left(e^{2\pi i h\cdot x}f\right)}(\xi)
	=(\tau_{h}\widehat f)(\xi).
\]

\begin{proof*}
	Analogamente ao item anterior,
	\begin{align*}
		\widehat{\left(e^{2\pi i h\cdot x}f\right)}(\xi)
		 & =\int_{\mathbb{R}^{n}}e^{2\pi i x\cdot h}f(x)e^{-2\pi i x\cdot\xi}\,dx \\
		 & =\int_{\mathbb{R}^{n}}f(x)e^{-2\pi i x\cdot(\xi-h)}\,dx                \\
		 & =\widehat f(\xi-h)
		=(\tau_{h}\widehat f)(\xi).
	\end{align*}
	\qedsymbol
\end{proof*}

\medskip
\noindent\textbf{(III)} Se $T:\mathbb{R}^{n}\to\mathbb{R}^{n}$ é um isomorfismo e
\[
	S\coloneqq(T^{*})^{-1},
\]
então
\[
	\widehat{(f\circ T)}(\xi)
	=|\det(T)|^{-1}(\widehat f\circ S)(\xi)
	=|\det(T)|^{-1}\widehat f(S\xi),
	\qquad \xi\in\mathbb{R}^{n}.
\]

\begin{proof*}
	Ainda pelo Teorema da Mudança de Variáveis, usando a definição do adjunto $T^{*}$ e a substituição $y=Tx$, temos
	\begin{align*}
		\widehat{(f\circ T)}(\xi)
		 & =\int_{\mathbb{R}^{n}}e^{-2\pi i x\cdot\xi}f(Tx)\,dx \\
		 & =\int_{\mathbb{R}^{n}}
		e^{-2\pi i (T^{-1}y)\cdot\xi}f(y)|\det(T^{-1})|\,dy     \\
		 & =|\det(T)|^{-1}
		\int_{\mathbb{R}^{n}}
		e^{-2\pi i y\cdot[(T^{-1})^{*}\xi]}f(y)\,dy             \\
		 & =|\det(T)|^{-1}\widehat f(S\xi),
	\end{align*}
	pois
	\[
		(T^{-1})^{*}=(T^{*})^{-1}=S.
	\]
	\qedsymbol
\end{proof*}

\begin{crl*}
	Dois casos particulares da propriedade acima merecem destaque.

	\medskip
	\noindent\textbf{(1º)} Se $T$ é uma rotação, isto é,
	\[
		T\circ T^{*}=I,
	\]
	então
	\[
		T^{-1}=T^{*}
		\quad\Longrightarrow\quad
		T=(T^{*})^{-1},
		\qquad |\det(T)|=1.
	\]
	Logo,
	\[
		\widehat{(f\circ T)}(\xi)=\widehat f(T\xi),
	\]
	isto é, $\mathcal{F}$ é invariante, ou comuta, por rotações.

	\medskip
	\noindent\textbf{(2º)} Se
	\[
		T_{t}x\coloneqq tx,
	\]
	é uma dilatação, com $t>0$, então
	\[
		T_{t}^{*}=T_{t},
		\qquad
		(T_{t}^{*})^{-1}=T_{t^{-1}},
		\qquad
		\det(T_{t})=t^{n}.
	\]
	Portanto,
	\[
		\widehat{(f\circ T_{t})}(\xi)
		=t^{-n}\widehat f(\xi/t)
		=(\widehat f)_{t}(\xi),
	\]
	na notação usada nas famílias regularizantes. Em particular,
	\[
		\boxed{\ \widehat{f(tx)}(\xi)=(\widehat f)_{t}(\xi)\ }.
	\]

	\medskip
	\noindent\textbf{(IV)} Vale
	\[
		\widehat{(f*g)}(\xi)=\widehat f(\xi)\widehat g(\xi),
		\qquad \forall \xi\in\mathbb{R}^{n}.
	\]
	Pela desigualdade de Young, se $f,g\in L^{1}$, então $f*g\in L^{1}$.
\end{crl*}
\begin{proof*}
	Um cálculo direto, usando o Teorema de Fubini, dá
	\begin{align*}
		\widehat{(f*g)}(\xi)
		 & =\int_{\mathbb{R}^{n}}e^{-2\pi i x\cdot\xi}(f*g)(x)\,dx \\
		 & =\int_{\mathbb{R}^{n}}e^{-2\pi i x\cdot\xi}
		\left[\int_{\mathbb{R}^{n}}f(x-y)g(y)\,dy\right]dx         \\
		 & =\int_{\mathbb{R}^{n}}g(y)
		\left[\int_{\mathbb{R}^{n}}e^{-2\pi i x\cdot\xi}f(x-y)\,dx\right]dy.
	\end{align*}
	Pelo item (I), o termo entre colchetes é
	\[
		e^{-2\pi i y\cdot\xi}\widehat f(\xi).
	\]
	Assim,
	\[
		\widehat{(f*g)}(\xi)
		=\widehat f(\xi)
		\int_{\mathbb{R}^{n}}g(y)e^{-2\pi i y\cdot\xi}\,dy
		=\widehat f(\xi)\widehat g(\xi).
	\]
	A passagem das integrais é permitida pelo Teorema de Fubini, pois
	\[
		(x,y)\longmapsto f(x-y)g(y)
	\]
	é integrável em $\mathbb{R}^{n}\times\mathbb{R}^{n}$ pelo Teorema de Tonelli, o que se prova usando que $f,g$ são Borel-mensuráveis. \qedsymbol
\end{proof*}

\medskip
\noindent\textbf{(V)} Se $f$ é tal que
\[
	x^{\alpha}f\in L^{1},
	\qquad \forall |\alpha|\leq m,
\]
então
\[
	\widehat f\in \mathcal{C}^{m}(\mathbb{R}^{n})
\]
e vale a igualdade
\[
	\partial_{\xi}^{\alpha}\widehat f(\xi)
	=\widehat{\left[(-2\pi i x)^{\alpha}f\right]}(\xi).
\]

\begin{tcolorbox}[
		skin=enhanced,
		title=Observação,
		fonttitle=\bfseries,
		colframe=black,
		colbacktitle=cyan!75!white,
		colback=cyan!15,
		colbacklower=black,
		coltitle=black,
		drop fuzzy shadow,
	]
	A igualdade significa que, se
	\[
		p_{\alpha}(x)=(-2\pi i x)^{\alpha}
	\]
	indica o polinômio correspondente, então
	\[
		\partial_{\xi}^{\alpha}\widehat f(\xi)
		=\widehat{(p_{\alpha}f)}(\xi).
	\]
	Assim, a Transformada de Fourier transforma multiplicação por polinômio em derivação: ``decaimento implica regularidade''.
\end{tcolorbox}

\begin{proof*}
	A prova é feita por indução a partir do caso $m=1$, derivando sob o sinal de integral. Uma vez que
	\[
		\frac{\partial}{\partial \xi_{j}}
		\left(e^{-2\pi i x\cdot\xi}f(x)\right)
		=-2\pi i x_{j}e^{-2\pi i x\cdot\xi}f(x),
	\]
	e que $x_{j}f\in L^{1}$, temos
	\begin{align*}
		\frac{\partial\widehat f}{\partial\xi_{j}}(\xi)
		 & =\frac{\partial}{\partial\xi_{j}}
		\int_{\mathbb{R}^{n}}e^{-2\pi i x\cdot\xi}f(x)\,dx                    \\
		 & =\int_{\mathbb{R}^{n}}
		\frac{\partial}{\partial\xi_{j}}
		\left(e^{-2\pi i x\cdot\xi}f(x)\right)dx                              \\
		 & =\int_{\mathbb{R}^{n}}(-2\pi i x_{j})e^{-2\pi i x\cdot\xi}f(x)\,dx \\
		 & =\widehat{\left[(-2\pi i x_{j})f\right]}(\xi).
	\end{align*}
	Como $(-2\pi i x_{j})f\in L^{1}$, segue da proposição anterior que
	\[
		\frac{\partial\widehat f}{\partial\xi_{j}}\in\mathcal{C}(\mathbb{R}^{n}).
	\]
	Sendo $j$ qualquer, $\widehat f\in\mathcal{C}^{1}(\mathbb{R}^{n})$. Repetindo o argumento, obtemos o caso geral. \qedsymbol
\end{proof*}

\medskip
\noindent\textbf{(VI)} Se, por outro lado, $f\in\mathcal{C}^{m}$ e satisfaz
\begin{itemize}
	\item[(i)] $\partial^{\alpha}f\in L^{1}$, para todo $|\alpha|\leq m$;
	\item[(ii)] $\partial^{\alpha}f\in\mathcal{C}_{0}(\mathbb{R}^{n})$, para $|\alpha|\leq m-1$, sendo
	      \[
		      \mathcal{C}_{0}(\mathbb{R}^{n})
		      =\{\text{funções contínuas que se anulam no infinito}\},
	      \]
\end{itemize}
então
\[
	\widehat{\partial^{\alpha}f}(\xi)
	=(2\pi i\xi)^{\alpha}\widehat f(\xi).
\]

\begin{tcolorbox}[
		skin=enhanced,
		title=Observação,
		fonttitle=\bfseries,
		colframe=black,
		colbacktitle=cyan!75!white,
		colback=cyan!15,
		colbacklower=black,
		coltitle=black,
		drop fuzzy shadow,
	]
	Novamente, a Transformada de Fourier transforma derivação em multiplicação por polinômio: ``regularidade implica decaimento''. Quando $n=1$, a condição (ii) é satisfeita automaticamente.
\end{tcolorbox}

\begin{proof*}
	A prova também se faz por indução a partir do caso $m=1$ e integrando por partes. Para uma derivada na variável $x_{j}$,
	\[
		\widehat{\left(\frac{\partial f}{\partial x_{j}}\right)}(\xi)
		=\int_{\mathbb{R}^{n}}e^{-2\pi i x\cdot\xi}\frac{\partial f}{\partial x_{j}}(x)\mathrm{d}x = \int_{\mathbb{R}^{n-1}}^{}\biggl[\int_{\mathbb{R}}^{}e^{-2\pi ix \xi }\frac{\partial^{}f(x)}{\partial x_{j}^{}} \mathrm{d}x_{j}\biggr] \mathrm{d}x_1 \dotsc \widehat{\mathrm{d}x_{j}}\dotsc \mathrm{d}x_{n}.
	\]
	Fixando as demais variáveis, a integral na coordenada $x_{j}$ calcula-se por partes:
	\begin{align*}
		\int_{\mathbb{R}}e^{-2\pi i x\cdot\xi}
		\frac{\partial f}{\partial x_{j}}(x)\,\mathrm{d}x_{j}
		 & =\lim_{N\to\infty}
		\left[
			e^{-2\pi i x\cdot\xi}f(x)
		\right]_{x_{j}=-N}^{x_{j}=N}                      \\
		 & \quad-\int_{-N}^{N}
		(-2\pi i\xi_{j})e^{-2\pi i x\cdot\xi}f(x)\,dx_{j} \\
		 & =(2\pi i\xi_{j})
		\int_{\mathbb{R}}e^{-2\pi i x\cdot\xi}f(x)\,dx_{j},
	\end{align*}
	pois o termo de bordo se anula, já que $f\in\mathcal{C}_{0}(\mathbb{R}^{n})$. Substituindo na integral nas variáveis restantes, obtemos
	\[
		\widehat{\left(\frac{\partial f}{\partial x_{j}}\right)}(\xi) = \int_{\mathbb{R}^{n-1}}^{}\biggl[\int_{\mathbb{R}}^{}e^{-2\pi ix \xi }f(x) \mathrm{d}x_{j}\biggr] \mathrm{d}x_1 \dotsc \widehat{\mathrm{d}x_{j}}\dotsc \mathrm{d}x_{n}
		=(2\pi i\xi_{j})\widehat f(\xi),
	\]
	completando o caso $m=1$. O caso geral segue por indução. \qedsymbol
\end{proof*}

\medskip
\noindent\textbf{(VII) Lema de Riemann--Lebesgue.} Para toda $f\in L^{1}(\mathbb{R}^{n})$, tem-se
\[
	\widehat f\in\mathcal{C}_{0}(\mathbb{R}^{n}),
\]
isto é,
\[
	\lim_{|\xi|\to\infty}\widehat f(\xi)=0.
\]

\begin{proof*}
	Isso é claro quando $f\in\mathcal{C}_{c}^{1}(\mathbb{R}^{n})$, porque, por (VI) e pela proposição inicial,
	\[
		|\xi_{j}\widehat f(\xi)|
		=\frac{1}{2\pi}
		\left|\widehat{\left(\partial_{j}f\right)}(\xi)\right|
		\leq\frac{1}{2\pi}\|\partial_{j}f\|_{L^{1}}.
	\]
	Logo,
	\[
		\left(\sum_{j=1}^{n}|\xi_{j}|\right)|\widehat f(\xi)|
		\leq\frac{1}{2\pi}\|\nabla f\|_{L^{1}}<\infty,
	\]
	e, portanto,
	\[
		\lim_{|\xi|\to\infty}|\widehat f(\xi)|=0.
	\]

	O caso geral resulta do fato de que $\mathcal{C}_{c}^{1}$ é denso em $L^{1}$, de que
	\[
		\mathcal{F}:L^{1}\longrightarrow\mathcal{C}_{B}
	\]
	é contínua e de que $\mathcal{C}_{0}$ é subespaço fechado de $\mathcal{C}_{B}$. Os detalhes são exercícios. Explicitamente, dado $\varepsilon>0$, escolhemos $f_{\varepsilon}\in\mathcal{C}_{c}^{1}$ com
	\[
		\|f_{\varepsilon}-f\|_{L^{1}}<\frac{\varepsilon}{2}.
	\]
	Então
	\[
		|\widehat f(\xi)|
		\leq
		|\widehat f(\xi)-\widehat f_{\varepsilon}(\xi)|
		+|\widehat f_{\varepsilon}(\xi)|.
	\]
	Pela continuidade de $\mathcal{F}$,
	\[
		|\widehat f(\xi)-\widehat f_{\varepsilon}(\xi)|
		\leq \|f-f_{\varepsilon}\|_{L^{1}}
		<\frac{\varepsilon}{2},
	\]
	e, como $\widehat f_{\varepsilon}\in\mathcal{C}_{0}$, existe $N\in\mathbb{N}$ tal que
	\[
		|\xi|\geq N
		\quad\Longrightarrow\quad
		|\widehat f_{\varepsilon}(\xi)|<\frac{\varepsilon}{2}.
	\]
	Daí $|\widehat f(\xi)|<\varepsilon$ para $|\xi|\geq N$. \qedsymbol
\end{proof*}

\begin{tcolorbox}[
		skin=enhanced,
		title=Observação,
		fonttitle=\bfseries,
		colframe=black,
		colbacktitle=cyan!75!white,
		colback=cyan!15,
		colbacklower=black,
		coltitle=black,
		drop fuzzy shadow,
	]
	Escrevendo
	\[
		e^{-2\pi i n x}=\cos(2\pi n x)+i\sin(2\pi n x),
	\]
	temos
	\[
		\widehat f(n)
		=\int\cos(2\pi n x)f(x)\,dx
		+i\int\sin(2\pi n x)f(x)\,dx.
	\]
	Quando $f$ é real, obtém-se o enunciado clássico do Lema de Riemann--Lebesgue (ou escrevendo $\xi=n/(2\pi)$, $n\in\mathbb{Z}$):
	\[
		\lim_{n\to\infty}\int\cos(2\pi n x)f(x)\,dx=0,
		\qquad
		\lim_{n\to\infty}\int\sin(2\pi n x)f(x)\,dx=0.
	\]
\end{tcolorbox}

\end{document}
